\begin{table}[H]
  \centering
  \caption{Concentración de la importancia SHAP por algoritmo y
  especificación, pobreza monetaria, 2013$\to$2016.}
  \label{tab:shap_variables_necesarias}
  \footnotesize
  \setlength{\tabcolsep}{5pt}
  \begin{tabular}{llccc}
    \toprule
    \textbf{Algoritmo} & \textbf{Espec.} & \textbf{\shortstack{Variables\\(total)}} & \textbf{\shortstack{\% SHAP en las\\10 primeras}} & \textbf{\shortstack{Variables para el\\80\% del SHAP}} \\
    \midrule
    HistGradientBoosting     & A & 169 & 69.3\% & 15 \\
    XGBoost                  & A & 169 & 58.7\% & 23 \\
    Logística regularizada   & A & 169 & 43.9\% & 28 \\
    Random Forest            & A & 169 & 59.8\% & 32 \\
    LightGBM                 & A & 169 & 44.4\% & 37 \\
    \addlinespace
    LightGBM                 & B & 165 & 71.6\% & 14 \\
    HistGradientBoosting     & B & 165 & 63.9\% & 18 \\
    XGBoost                  & B & 165 & 54.7\% & 26 \\
    Logística regularizada   & B & 165 & 45.8\% & 28 \\
    Random Forest            & B & 165 & 54.9\% & 34 \\
    \bottomrule
  \end{tabular}
  \begin{minipage}{0.95\textwidth}
    \vspace{4pt}
    \footnotesize \textit{Nota:} Importancia SHAP calculada sobre el conjunto
    de prueba. Variables originales: las \emph{dummies one-hot} y los
    indicadores de faltante se suman a su variable de origen antes de
    calcular la masa acumulada. Espec.: especificación (A con ingreso y
    gasto del hogar, B sin ellos). Fuente: cálculos propios.
  \end{minipage}
\end{table}